\documentclass[11pt,a4paper]{article}
\usepackage[margin=2.2cm]{geometry}
\usepackage{amsmath,amssymb,bm}
\usepackage[table]{xcolor}
\usepackage{booktabs,tabularx,array}
\usepackage{enumitem}
\usepackage[most]{tcolorbox}
\usepackage{tikz}
\usetikzlibrary{arrows.meta,positioning,calc,angles,quotes}
\usepackage{fontspec}
\usepackage[colorlinks=true,linkcolor=blue!60!black,urlcolor=blue!60!black]{hyperref}
\usepackage{polyglossia}
\setmainlanguage[numerals=maghrib]{arabic}
\setotherlanguage{english}
\newfontfamily\arabicfont[Script=Arabic,Scale=1.08]{Amiri}
\newfontfamily\englishfont{Latin Modern Roman}
\setmonofont{DejaVu Sans Mono}[Scale=0.82]
\definecolor{cblue}{HTML}{2B6CB0}\definecolor{cgreen}{HTML}{2F855A}\definecolor{cred}{HTML}{C53030}
\definecolor{corange}{HTML}{DD6B20}\definecolor{cgray}{HTML}{718096}
\newtcolorbox{idea}[1][]{colback=cblue!5,colframe=cblue,title={#1},fonttitle=\bfseries,breakable}
\newtcolorbox{why}[1][]{colback=cgreen!5,colframe=cgreen,title={#1},fonttitle=\bfseries,breakable}
\newtcolorbox{warn}[1][]{colback=cred!5,colframe=cred,title={#1},fonttitle=\bfseries,breakable}
\newtcolorbox{example}[1][]{colback=corange!5,colframe=corange,title={#1},fonttitle=\bfseries,breakable}
\newtcolorbox{codebox}[1][]{colback=cgray!7,colframe=cgray,title={#1},fonttitle=\bfseries,breakable}
\newcommand{\en}[1]{\textenglish{#1}}
\newcommand{\code}[1]{\LR{\texttt{#1}}}
\newcommand{\vv}{\bm v}\newcommand{\pp}{\bm p}\newcommand{\uu}{\bm u}
\setlist{itemsep=2pt,topsep=3pt}
\begin{document}

\begin{center}
{\LARGE\bfseries شرح الجزء الأول من التشغيل المؤجَّل}\\[6pt]
{\large إعادة تشغيل المقارنات الأساسية بعد التدقيق البحثي، وبيانات الفيزياء الواقعية}\\[6pt]
{\color{cgray} نسخة عربية موسّعة، كُتبت في \LR{2026-11-02} للفصل الثاني من الأطروحة (تقدير حالة الكرة بالرؤية).\\
كل رقم ومعادلة هنا مطابق للكود في المستودع. المصطلحات الإنجليزية تُذكر بين قوسين لتسهيل الربط بالكود والأوراق العلمية.}
\end{center}

\tableofcontents
\newpage

%=====================================================================================================
\section{الجزء الأول في صفحة واحدة}

الجزء الأول هو النصف الأول من الملف \code{experiments/run\_overnight.sh}. يتكوّن من سبع خطوات ويستغرق نحو ثلاث ساعات
ونصف (الضبط وحده استغرق فعلياً نحو $100$ دقيقة). الخطوات من $1$ إلى $5$ تعيد حساب نتائج المرحلة صفر بعد تطبيق كل
إصلاحات التدقيق، والخطوتان $6$ و$7$ تفحصان وتُصيِّران (\en{render}) بيانات الفيزياء الواقعية الجديدة.

\begin{LTR}
\begin{center}
\begin{tikzpicture}[node distance=3.5mm, box/.style={draw,rounded corners,fill=cblue!8,minimum width=12.5cm,align=right,inner sep=4pt}]
\node[box] (a) {\textarabic{\textbf{١. الملاءمة}: ملاءمة الاحتكاك الثابت $c$ ومعامل الارتداد $E$ على رميات \emph{التدريب} \hfill $\sim15$ دقيقة}};
\node[box,below=of a] (b) {\textarabic{\textbf{٢. الضبط}: ضبط الضوضاء $(\sigma,q)$ لكل عائلة مرشحات على رميات \emph{التحقق} \hfill $\sim100$ دقيقة}};
\node[box,below=of b] (c) {\textarabic{\textbf{٣. المقارنات الأساسية}: $13$ مقدِّراً $\times$ $11$ حالة حساس على التحقق والاختبار التطويري \hfill $\sim17$ دقيقة}};
\node[box,below=of c] (d) {\textarabic{\textbf{٤. الاختبارات المقترنة}: هل الطريقة \en{A} أفضل فعلاً من \en{B}؟ \hfill ثوانٍ}};
\node[box,below=of d] (e) {\textarabic{\textbf{٥. $\mu$ مقابل $E$}: إعادة فحص الفجوة على التدريب والتحقق فقط \hfill $\sim2$ دقيقة}};
\node[box,below=of e,fill=cgreen!10] (f) {\textarabic{\textbf{٦. الفحص}: فحص الفيزياء الواقعية مقابل معادلات الكتب و$500$ رمية \hfill $\sim3$ دقائق}};
\node[box,below=of f,fill=cgreen!10] (g) {\textarabic{\textbf{٧. التصيير}: تمرير $500$ رمية واقعية عبر الكاميرا (أُجِّل بطلبك) \hfill $\sim90$ دقيقة}};
\foreach \x/\y in {a/b,b/c,c/d,d/e,e/f,f/g} \draw[-{Stealth},thick] (\x) -- (\y);
\end{tikzpicture}
\end{center}
\end{LTR}

\begin{why}[لماذا نحتاج الجزء الأول أصلاً؟]
وجد التدقيق البحثي (\code{docs/research\_audit\_2026-11-01.md}) مشكلات قد تؤدي إلى رفض البحث عند النشر. الجزء الأول هو المكان
الذي تتحول فيه الإصلاحات إلى أرقام:
\begin{itemize}
\item \textbf{\en{A1}} (الجريمة العكسية \en{inverse crime}): كان المحاكي يستخدم معادلات المقدِّر نفسها حرفياً. الحل مجموعة بيانات ثانية
بفيزياء واقعية (الخطوتان $6$ و$7$).
\item \textbf{\en{A2}}: نظرنا سابقاً إلى مجموعة الاختبار. من الآن تُعرض النتائج على التحقق وعلى «اختبار تطويري»، ومجموعة
الاختبار النهائية الحقيقية لن تُصيَّر إلا في النهاية.
\item \textbf{\en{A3}}: لم تكن هناك فواصل ثقة، وكانت بضع رميات تسيطر على النتيجة. الحل إحصاء لكل رمية، مع فواصل
\en{bootstrap} واختبارات مقترنة (الخطوتان $3$ و$4$).
\item \textbf{\en{B1}} و\textbf{\en{B2}}: إصلاحان في القياس: تصحيح العمق على طول الشعاع، وطريقة قياس جديدة بتقاطع الشعاع مع
مستوى الأرض.
\item \textbf{\en{B3}}: لم تكن المقارنات عادلة. الثوابت تُلاءَم الآن على التدريب، وكل عائلة مرشحات تُضبط وحدها (الخطوتان $1$ و$2$).
\item \textbf{\en{B4}}: كان كاشف الإطارات المجمَّدة ناجحاً «بالبناء» لا بالجدارة. حالة حساس جديدة تجعله يفشل عمداً.
\end{itemize}
\end{why}

يشرح باقي هذا المستند كل ما تحتاجه: العالم المحاكى، والكاميرا، والمرشحات، والمقاييس، والإحصاء، ثم كل خطوة بالتفصيل
الرياضي، مع توسيع خاص للطرق الجديدة. وفي النهاية ملخص للنتائج الفعلية للتشغيل.

%=====================================================================================================
\section{الخلفية الأولى: العالم المحاكى (الفيزياء «المطابقة»)}

\subsection{الساحة والكرة والذراع}
\begin{itemize}
\item أرضية مساحتها $2.5\times2.5$ م، والجدران الداخلية عند $x,y=\pm1.25$ م.
\item الكرة نصف قطرها $r=0.03$ م، وتتحرك في المستوى فقط (مفصلا انزلاق، بلا دوران ولا طيران).
\item حدّ الجدار لمركز الكرة: $L=1.25-0.03=1.22$ م.
\item الذراع روبوت \en{Franka Panda} مجمَّد في وضعه الأساسي في المركز. الجزء الوحيد على ارتفاع الكرة هو القاعدة، وتُمثَّل
لمركز الكرة بصندوق مركزه $(c_x,c_y)=(-0.041,0)$ وأنصاف أبعاده $(h_x,h_y)=(0.143,0.125)$ م (القاعدة الحقيقية مضافاً إليها نصف قطر الكرة).
\item خطوة المحاكاة الزمنية: $\Delta t=2$ مللي ثانية.
\end{itemize}

\subsection{معادلات الحركة (الملف \code{arena\_sim.py}، الصنف \code{Arena})}
برنامج \en{MuJoCo} هنا يقوم بالتكامل الزمني والتصيير فقط؛ كل التلامسات معطَّلة، وفيزياء الكرة مكتوبة بلغة بايثون:
\begin{enumerate}
\item \textbf{الاحتكاك}: تباطؤ ثابت $\mu$ (م/ث$^2$) عكس اتجاه السرعة:
\[ \dot{\vv}=-\mu\,\frac{\vv}{|\vv|}. \]
إذا صارت $|\vv|\le\max(10^{-3},\mu\Delta t)$ توقفت الكرة تماماً، لأن الاحتكاك لا يجوز أن يعكس اتجاه الحركة.
\item \textbf{التكامل}: أويلر شبه الضمني: أولاً $\vv\leftarrow\vv+\dot{\vv}\Delta t$ ثم $\pp\leftarrow\pp+\vv\Delta t$.
\item \textbf{الارتداد عن الجدار}: إذا تجاوز المركز $x>L$ وهو يتحرك للخارج، يصبح $x\leftarrow L$ و$v_x\leftarrow -E\,v_x$،
وتبقى السرعة المماسية كما هي.
\item \textbf{الارتداد عن الذراع}: إذا دخل المركز الصندوق يُدفع للخارج عبر الوجه ذي الاختراق الأصغر، وتصبح السرعة العمودية
$-E\,v_n$ إن كانت متجهة للداخل.
\end{enumerate}
في كل رمية يُسحب $\mu\sim U(0.2,1.5)$ م/ث$^2$ و$E\sim U(0.5,0.95)$ مرة واحدة ويبقيان ثابتين طوال الرمية.

\begin{warn}[تنبيه مهم: $\mu$ هنا تباطؤ وليس معامل احتكاك]
في هذا المشروع وحدة $\mu$ هي م/ث$^2$. المعامل اللابعدي هو $\mu/g$؛ فمثلاً $\mu=0.85$ م/ث$^2$ تعادل معاملاً قدره $0.087$.
يجب أن تذكر الأطروحة هذا بوضوح (بند التدقيق \en{C1}).
\end{warn}

\subsection{الرميات الخمسمئة والتقسيم}
خمسون سرعة ($0.1,0.2,\dots,5.0$ م/ث)، لكل منها $10$ رميات، أي $500$ رمية، بنقطة بداية عشوائية (لا تكون داخل الذراع أو قربه)
واتجاه عشوائي (البذرة $42$)، ومدة كل رمية $10$ ثوانٍ.
التقسيم (\code{split\_500.json}، البذرة $2026$): لكل سرعة تذهب $7$ رميات إلى \textbf{التدريب}، والثلاث الباقية إلى التحقق
والاختبار بنسبة $2+1$ أو $1+2$ بالتناوب، فالمجموع $350$ تدريب و$75$ تحقق و$75$ اختبار.

\begin{idea}[الأدوار الثلاثة للبيانات]
\begin{itemize}
\item \textbf{التدريب}: بيانات نتعلّم أو نلائم عليها أرقاماً (الخطوة $1$، ولاحقاً \en{ekf-$\mu$}).
\item \textbf{التحقق}: بيانات نختار بها الإعدادات (الخطوة $2$).
\item \textbf{الاختبار}: بيانات للعرض فقط. إذا اخترت أي شيء بالنظر إلى نتائج الاختبار، لم يعد الاختبار تقديراً صادقاً
لأداء الطريقة على بيانات جديدة.
\end{itemize}
لأننا نظرنا إلى مجموعة الاختبار في المرحلة صفر، صار اسمها «الاختبار التطويري» (\en{test (dev)}). أما مجموعة الاختبار
النهائية فستكون رميات جديدة بالبذرة $20261101$ (الملف \code{final\_test\_plan.json})، تُصيَّر مرة واحدة بعد تجميد كل الطرق،
وتُقيَّم مرة واحدة، وتُعرض كل أرقامها.
\end{idea}

%=====================================================================================================
\section{الخلفية الثانية: الكاميرا والقياس}

\subsection{كاميرا الثقب (\en{pinhole})}
الكاميرا تنظر عمودياً إلى الأسفل من ارتفاع $H=3$ م، والصورة $480\times480$ بكسل، ومجال الرؤية الرأسي $49^\circ$.
البعد البؤري بالبكسل:
\[ f=\frac{480/2}{\tan(49^\circ/2)}=\frac{240}{\tan 24.5^\circ}=526.6\ \text{px}. \]
مركز الصورة $(c_u,c_v)=(239.5,\,239.5)$. النقطة التي عمقها $Z$ (المسافة على المحور البصري) وتظهر في البكسل $(u,v)$ موقعها:
\[ x=\frac{u-c_u}{f}\,Z,\qquad y=-\frac{v-c_v}{f}\,Z \]
(الإشارة السالبة لأن صفوف الصورة تزداد نحو الأسفل).

\begin{example}[حجم البكسل الواحد على الأرض]
عند عمق مركز الكرة $Z=3-0.03=2.97$ م يغطي البكسل الواحد $2.97/526.6=5.6$ مم. عرض المشهد كله $2.73$ م، فتدخل الساحة ومعها هامش
نحو $5$ سم؛ وهذه هي القاعدة التي اختير بها مجال الرؤية $49^\circ$.
\end{example}

\subsection{الكشف عن الكرة}
يعطي المُصيِّر صورة ألوان وصورة عمق. الكرة خضراء: البكسل ينتمي للكرة إذا زادت قيمة الأخضر على الأحمر وعلى الأزرق بأكثر من
$40$. نحتاج على الأقل $3$ بكسلات. \textbf{المركز الهندسي} $(u,v)$ هو متوسط إحداثياتها، ويُقرأ العمق $z_s$ عند أقرب بكسل
إليه، وهو عمق \emph{سطح} الكرة لا مركزها.
إذا غطّت الذراع الكرة كلياً فلا قياس. وإذا غطّت جزءاً منها انزاح مركز الجزء المرئي، وهذا يحدث حتى في الحالة
«المثالية» ويسبب أخطاء تصل إلى نحو $7$ سم (بند \en{B2}).

\subsection{طريقة جديدة \#1: تصحيح العمق على طول الشعاع (بند \en{B2})}
\textbf{المشكلة.} الكاميرا ترى سطح الكرة، ونحن نريد مركزها. الكود القديم كان يضيف نصف القطر كاملاً: $Z_c=z_s+r$. هذا صحيح تماماً
فقط للكرة الواقعة مباشرة تحت الكاميرا.

\textbf{الهندسة.} الشعاع من الكاميرا إلى مركز الصورة الكروية يمر (تقريباً تماماً) بمركز الكرة. على هذا الشعاع يكون السطح
أقرب إلى الكاميرا من المركز بمسافة $r$ بالضبط. لكن العمق $Z$ يُقاس على المحور البصري لا على الشعاع. إذا صنع الشعاع زاوية
$\theta$ مع المحور البصري، فإن المسافة $r$ على الشعاع تقابل $r\cos\theta$ فقط على المحور:

\begin{LTR}
\begin{center}
\begin{tikzpicture}[scale=1.25]
\coordinate (C) at (0,5); \coordinate (B) at (3.4,1.2);
\draw[thick] (-0.35,5) -- (0.35,5); \node[above] at (C) {\textarabic{الكاميرا}};
\draw[dashed,gray] (C) -- (0,0.3) node[below] {\textarabic{المحور البصري}};
\draw[thick,cblue] (C) -- ($(C)!1.15!(B)$);
\draw[fill=cgreen!25] (B) circle (0.75);
\fill (B) circle (1.5pt); \node[right=4pt] at (B) {\textarabic{المركز}};
\coordinate (S) at ($(B)!0.75cm!(C)$);
\fill[cred] (S) circle (1.8pt);
\draw[cred,-{Stealth}] ($(S)+(-1.4,0.9)$) node[above left,text=cred] {\textarabic{نقطة السطح (يقيسها العمق)}} -- ($(S)+(-0.08,0.06)$);
\draw[dotted,thick] (S) -- (-1.0,0 |- S); \draw[dotted,thick] (B) -- (-1.0,0 |- B);
\draw[<->,cred,thick] (-1.0,0 |- S) -- (-1.0,0 |- B) node[midway,left] {$r\cos\theta$};
\draw[<->] ($(S)!0!(B)+(0.35,0.3)$) -- ($(B)+(0.35,0.3)$) node[midway,above right=-1pt] {$r$};
\draw (0,4.0) arc (-90:-50:1.0); \node at (0.3,3.65) {$\theta$};
\end{tikzpicture}
\end{center}
\end{LTR}

\[ Z_c=z_s+r\cos\theta,\qquad \cos\theta=\frac{1}{\sqrt{1+a^2+b^2}},\quad a=\frac{u-c_u}{f},\ b=\frac{v-c_v}{f}. \]

\textbf{لماذا $\cos\theta=1/\sqrt{1+a^2+b^2}$؟} اتجاه الشعاع في إحداثيات الكاميرا هو $(a,b,1)$ (نتحرك $a$ أفقياً و$b$ رأسياً لكل
وحدة عمق). جيب تمام الزاوية بين هذا المتجه والمحور $(0,0,1)$ هو المركبة الثالثة مقسومة على الطول: $1/\sqrt{a^2+b^2+1}$.

\begin{example}[حجم الخطأ القديم]
على بعد $1.5$ م من المركز: $\cos\theta=2.97/\sqrt{2.97^2+1.5^2}=0.893$. خطأ العمق القديم $r(1-\cos\theta)=0.03\times0.107=3.2$ مم،
وهو يعطي خطأ موقع قدره $3.2\times1.5/2.97=1.6$ مم نحو الخارج. على البيانات الفعلية انخفض الانحياز من $+3.0$ إلى $+1.6$ مم
في النطاق $1.3$ إلى $1.8$ م.

الباقي ($+1.6$ مم) هو مقياس صافٍ قدره $1.0008$ (أي $0.08\%$)، ويظهر حتى في طريقة المستوى التي لا تستخدم العمق؛ إذن مصدره المركز
الهندسي في الصورة لا العمق. هو أقل من ثلث بكسل، فوُثِّق كخطأ متبقٍّ يشبه خطأ المعايرة.
\end{example}

\subsection{طريقة جديدة \#2: تقاطع الشعاع مع مستوى الأرض (بند \en{B1})}
\textbf{الفكرة.} مركز الكرة يقع دائماً على مستوى معروف هو $z=r$ (الكرة على الأرض). إذن لا نحتاج قياس العمق إطلاقاً: يكفي أن نقطع
الشعاع بهذا المستوى. المسافة من الكاميرا إلى المستوى على المحور البصري معروفة: $H-r=2.97$ م، فيكون:
\[ x=\frac{u-c_u}{f}\,(H-r),\qquad y=-\frac{v-c_v}{f}\,(H-r). \]

\textbf{تحليل الأخطاء: لماذا يُتوقع أن تكون أفضل؟}
\begin{itemize}
\item بطريقة العمق، خطأ عمق $\delta z$ يحرّك الموقع بمقدار $\delta x=x\,\delta z/z$. مع ضوضاء شبيهة بكاميرا \en{D435}
($\sigma_Z=33$ مم) وعند $x=1.2$ م: $1.2\times33/2.97=13$ مم. عند مركز الصورة يكون قريباً من الصفر، أي أن الخطأ يكبر نحو الأطراف.
\item بطريقة المستوى لا يبقى إلا خطأ البكسل: $0.3$ بكسل $\times5.6$ مم $=1.7$ مم، وهو تقريباً متساوٍ في كل مكان.
\item الثمن: يجب معرفة ارتفاع الكاميرا. لو كان الارتفاع المفترض خاطئاً بمقدار $\delta H$ لصار كل موقع مضروباً في
$(H-r+\delta H)/(H-r)$، أي خطأ مقياس. فخطأ $1$ سم في الارتفاع يعطي $1.2\times0.01/2.97=4$ مم عند الحافة، وهو انحياز ثابت لا ضوضاء.
لهذا يوجد في الكود المعامل \code{plane\_dz} لاختبار هذه الحساسية لاحقاً.
\item ميزة إضافية: لا يتأثر بالعمق الخاطئ عند الإخفاء الجزئي بالذراع (حين يقع بكسل العمق على الذراع نفسها). أقصى خطأ في الحالة
المثالية انخفض من $67$ إلى $28$ مم.
\end{itemize}
هذا أول ما سيقترحه مُحكِّم الورقة، فأضفنا حالتين: \code{plane} و\code{all\_plane}.

\subsection{آثار الحساس (الملف \code{sensing/sensor.py})}
تُطبَّق هذه الآثار \emph{بعد} التصيير على الكشوفات المثالية المخزَّنة، فيمكن رؤية الرمية نفسها عبر كل حالات الحساس.
\begin{enumerate}
\item \textbf{ضوضاء العمق}: $z_s\leftarrow z_s+\mathcal N(0,\sigma_Z^2)$ حيث $\sigma_Z=k\,Z^2$، وهذا الشكل مأخوذ من كاميرات
الاستريو (خطأ العمق ينمو مع مربع المسافة). القيمة $k=8.2\times10^{-4}$ تعطي $\sigma_Z=7.2$ مم (كاميرا زمن طيران جيدة، حالة
«\en{noise}»)، و$k=3.74\times10^{-3}$ تعطي $33$ مم (مثل \en{Intel D435} على $3$ م، حالة «\en{noise\_real}»).
\item \textbf{اهتزاز البكسل}: $u,v\leftarrow u,v+\mathcal N(0,0.3^2)$ بكسل، أي $1.7$ مم على الأرض.
\item \textbf{فقدان الإطارات}: يحدث على دفعات وفق سلسلة ماركوف ذات حالتين (انظر الصندوق أدناه).
\item \textbf{أنماط الفقدان}: \en{flagged} (يُبلَّغ المرشح بعدم وجود قياس)؛ \en{silent\_hold} (يعيد الحساس آخر قيمة بتطابق تام
بِتّاً ببت ويدّعي أنها صالحة)؛ و\textbf{الجديد} \en{silent\_noisy} (يعيد آخر كشف لكن بضوضاء \emph{جديدة}، فلا يكون متطابقاً).
\item \textbf{التأخير}: الإطار الملتقط في اللحظة $t$ يصل في $t+0.04$ ث، والمرشح يعرف زمن الالتقاط.
\end{enumerate}

\begin{idea}[سلسلة ماركوف للفقدان، مع الرياضيات]
لتكن $q$ احتمال الانتقال في كل إطار من \emph{مفقود} إلى \emph{مرئي}، و$p$ احتمال الانتقال من \emph{مرئي} إلى \emph{مفقود}.
\begin{itemize}
\item طول دفعة الفقدان له توزيع هندسي متوسطه $1/q$ إطاراً، فنختار $q=1/g$ مع متوسط فجوة $g=6$ إطارات.
\item في التوازن (على المدى الطويل) يساوي التدفق من مرئي إلى مفقود التدفق المعاكس: $\pi_{\text{seen}}\,p=\pi_{\text{lost}}\,q$،
ومع $\pi_{\text{seen}}+\pi_{\text{lost}}=1$ نحصل على $\pi_{\text{lost}}=p/(p+q)$.
\item نريد أن تساوي هذه النسبة معدل الفقدان $r$:
\[ \frac{p}{p+q}=r\ \Rightarrow\ p=\frac{r\,q}{1-r}. \]
\item مثال: $r=0.05$ و$g=6$ تعطي $q=0.167$ و$p=0.00877$، والتحقق: $p/(p+q)=0.05$.
\end{itemize}
لسلسلة الفقدان مولِّد عشوائي خاص بها، فكل الحالات ذات معاملات الفقدان نفسها تفقد \emph{الإطارات نفسها} بالضبط. هذا يجعل
المقارنة بين الحالات «مقترنة».
\end{idea}

\subsection{طريقة جديدة \#3: الإطارات المجمَّدة غير القابلة للكشف (بند \en{B4})}
\textbf{المشكلة.} في النمط \en{silent\_hold} يكرر الحساس القيمة نفسها بتطابق تام، والمُغلِّف \en{StaleGuard} يعتبر أي قياس مطابق
تماماً للسابق «مفقوداً». لذلك كان الكاشف ناجحاً بنسبة $100\%$ تلقائياً: أعطى الضبط في المرحلة صفر الرقم نفسه ($61.5$ مم)
للفقدان المُبلَّغ والفقدان الصامت. سيقول المحكِّم: المشكلة «محلولة بالبناء».

\textbf{الحل.} النمط \en{silent\_noisy}: الإطار المفقود يعيد آخر كشف $(u,v,z)$ لكن تُضاف إليه ضوضاء جديدة (اهتزاز وعمق)،
فيبدو قياساً عادياً صالحاً. يحدث هذا في الواقع مثلاً حين يعيد خط المعالجة معالجة صورة قديمة بعيِّنات ضوضاء جديدة.
الحالة الجديدة \code{dropout\_silent\_noisy} تقيس كم يكلّف ذلك، ومقارنتها مع حالة \en{noise} (الضوضاء نفسها بلا فقدان) هي
الإجابة.

\textbf{ملاحظة تحقق مهمة:} تأكدنا أن الكرة المتحركة لا تعطي أبداً كشفين مثاليين متطابقين ($0$ من $78\,044$)، أي أن
\en{StaleGuard} لا يرمي قياسات صحيحة في الحالة المثالية.

\subsection{حالات الحساس الإحدى عشرة}
\begin{center}\small
\begin{tabularx}{\linewidth}{@{}l>{\raggedleft\arraybackslash}X>{\raggedleft\arraybackslash}X@{}}\toprule
الاسم & ما تضيفه & السؤال الذي تجيب عنه\\\midrule
\en{ideal} & لا شيء (لكن: انحياز المركز والإخفاء، \en{B2}) & أفضل حالة\\
\en{noise} & $\sigma_Z=7$ مم، اهتزاز $0.3$ بكسل & كاميرا عمق جيدة\\
\en{dropout\_flagged} & فقدان $5\%$ بدفعات طولها $6$، مُبلَّغ & الإطارات المفقودة\\
\en{dropout\_silent} & الفقدان نفسه، تكرار متطابق & كاميرا مجمَّدة\\
\en{delay} & تأخير $40$ مللي ثانية & الكمون\\
\en{all} & ضوضاء + فقدان مُبلَّغ + تأخير & كل شيء (متفائل)\\
\en{noise\_real} & $\sigma_Z=33$ مم، اهتزاز $0.3$ بكسل & عمق شبيه بـ \en{D435}\\
\en{all\_real} & \en{noise\_real} + فقدان + تأخير & كل شيء (واقعي)\\\midrule
\textbf{\en{plane}} (جديدة) & مثل \en{noise\_real} لكن الموقع من المستوى بلا عمق & \en{B1}: هل نحتاج العمق؟\\
\textbf{\en{all\_plane}} (جديدة) & مثل \en{all\_real} لكن بقياس المستوى & \en{B1} في الحالة الصعبة\\
\textbf{\en{dropout\_silent\_noisy}} (جديدة) & ضوضاء + فقدان يُكرَّر بضوضاء جديدة & \en{B4}: إطارات قديمة لا تُكشف\\\bottomrule
\end{tabularx}\end{center}

\begin{idea}[الاقتران داخل حالتي المستوى]
في نمط المستوى ما زال الكود يسحب رقم ضوضاء العمق العشوائي ثم يرميه. السبب: لو لم يسحبه لتغيّر تسلسل الأرقام العشوائية
التالية. بهذه الحيلة تتلقى \en{plane} قيم اهتزاز البكسل نفسها تماماً التي تتلقاها \en{noise\_real} بالبذرة نفسها، فلا تختلف
الحالتان إلا في طريقة حساب $(x,y)$، والمقارنة المقترنة تقيس هذا الفرق وحده.
\end{idea}

%=====================================================================================================
\section{الخلفية الثالثة: المقدِّرات}

كل المقدِّرات لها الواجهة نفسها: \code{update(t, xy, valid)} في كل إطار، و\code{forecast(dt)} التي تتنبأ بالحالة بعد
\code{dt} ثانية من زمن المرشح.

\subsection{\en{fd}: الفروق المحدودة}
الموقع هو آخر قياس، والسرعة هي فرق آخر قياسين صالحين مقسوماً على فرق زمنيهما، والتنبؤ خط مستقيم: $\pp(t+\Delta)=\pp+\vv\Delta$.
لا احتكاك ولا جدران: أبسط مقارنة ممكنة.

\subsection{نموذج الحركة في مرشحات \en{EKF} (\code{estimators/model.py})}
الفيزياء نفسها التي في المحاكي المطابق، مطبَّقة على خطوات فرعية $h\le5$ مللي ثانية أثناء الترشيح و$h\le10$ للتنبؤ:
\begin{enumerate}
\item الاحتكاك: إذا $c\,h\ge|\vv|$ توقفت الكرة، وإلا $\vv\leftarrow\vv-c\,h\,\vv/|\vv|$.
\item الحركة: $\pp\leftarrow\pp+\vv h$.
\item الجدران: إذا $x>L$ يصبح $x\leftarrow L-E(x-L)$ و$v_x\leftarrow -Ev_x$. التجاوز يُعكَس ويُقصَّر بالعامل $E$، وهذا هو الموقع
الصحيح لو حدث الارتداد داخل الخطوة.
\item صندوق الذراع: اختبار «مسح» على القطعة المستقيمة من الموقع القديم إلى الجديد يحدد أي وجه عُبر أولاً ثم يعكس بالطريقة نفسها.
لأن القطعة كلها تُفحص، لا تستطيع الخطوة الطويلة أن تختار وجهاً خاطئاً أو أن تقفز فوق زاوية.
\end{enumerate}

\subsection{مرشح كالمان الموسَّع (\en{EKF19}، عن \en{Liu} وزملائه $2021$)}
\textbf{الحالة} $\bm s=(x,y,v_x,v_y)$. \textbf{عدم اليقين}: مصفوفة التغاير $P$ ($4\times4$)؛ القطر يحمل التباينات، وخارج القطر
يحمل الروابط بين الأرقام (مثلاً: إذا كان الموقع أبعد مما أظن، فالسرعة على الأرجح أكبر).

\paragraph{البداية.} عند أول قياس صالح $\bm z$: $\bm s=(z_x,z_y,0,0)$ و$P_0=\mathrm{diag}(\sigma^2,\sigma^2,\sigma_{v0}^2,\sigma_{v0}^2)$
مع $\sigma_{v0}=3$ م/ث، لأن الإطار الأول لا يقول شيئاً عن السرعة.

\paragraph{التنبؤ} (على خطوات فرعية $h$ حتى زمن القياس التالي):
\[ \bm s\leftarrow g(\bm s),\qquad P\leftarrow FPF^\top+Q, \]
حيث $g$ نموذج الحركة و$F=\partial g/\partial\bm s$ مصفوفة جاكوبي (مصفوفة المشتقات). للاحتكاك:
\[ \frac{\partial\vv'}{\partial\vv}=I-\frac{ch}{|\vv|}\left(I-\hat\vv\hat\vv^\top\right),\qquad \frac{\partial\pp'}{\partial\vv}=hI. \]
معناها: الاحتكاك يقلّص عدم اليقين \emph{عرضياً} على اتجاه الحركة قليلاً، لا \emph{على طوله}، لأن الحدّ $\hat\vv\hat\vv^\top$
يحذف المركبة الطولية.

\begin{warn}[تقريب معروف في هذا المرشح الأساسي (بند \en{B5})]
لا يُدخل \en{EKF19} الارتداد في $F$: بعد الارتداد يُعكس المتوسط ولا تُعكس $P$. المرشح \en{ekf-$\mu$} في الجزء الثاني يفعل ذلك
بشكل صحيح. في المقارنة الأساسية نُبقي نسخة \en{Liu} عن قصد ونسمّيها بوضوح.
\end{warn}

\paragraph{اشتقاق ضوضاء العملية $Q$.}
$Q$ تقول كم يمكن أن تختلف الحركة الحقيقية عن النموذج في خطوة واحدة. الاختيار المعياري «تسارع ضوضاء بيضاء»: تتلقى الكرة
تسارعات عشوائية $a(t)$ كثافتها الطيفية $q$ (الوحدة م$^2$/ث$^3$). خلال خطوة $h$ تغيّر السرعة بمقدار $\delta v=\int_0^h a\,dt$ والموقع
بمقدار $\delta p=\int_0^h(h-t)\,a\,dt$. ومع $\mathbb E[a(t)a(t')]=q\,\delta(t-t')$:
\[ \mathbb E[\delta v^2]=qh,\qquad \mathbb E[\delta p\,\delta v]=q\frac{h^2}{2},\qquad \mathbb E[\delta p^2]=q\frac{h^3}{3}. \]
فلكل محور $Q=q\begin{pmatrix}h^3/3&h^2/2\\h^2/2&h\end{pmatrix}$، وهذا بالضبط ما يبنيه الكود. $q$ كبيرة تعني «أثق بنموذج حركتي قليلاً».

\paragraph{التصحيح (التحديث)} بقياس الموقع $\bm z$، و$H=[I_2\ 0]$:
\begin{align*}
\bm y&=\bm z-H\bm s &&\text{الابتكار: المقاس ناقص المتوقع}\\
S&=HPH^\top+R,\quad R=\sigma^2I_2 &&\text{الحجم المتوقع للابتكار}\\
K&=PH^\top S^{-1} &&\text{كسب كالمان}\\
\bm s&\leftarrow\bm s+K\bm y,\qquad P\leftarrow(I-KH)P.
\end{align*}

\begin{example}[الكسب في بُعد واحد]
إذا كان التنبؤ غير متيقن بمقدار $10$ مم والكاميرا بمقدار $5$ مم، فإن $K=10^2/(10^2+5^2)=0.8$: يتحرك المرشح $80\%$ من الطريق نحو
القياس. مع كاميرا دقتها $1$ مم: $K=100/101\approx0.99$. ومع كاميرا خطؤها $30$ مم: $K=100/1000=0.1$.
\end{example}

\paragraph{البوابة والسياسات.} $d^2=\bm y^\top S^{-1}\bm y$ يقيس مدى «مفاجأة» القياس. إذا كان المرشح صحيحاً يتبع $d^2$ توزيع
$\chi^2$ بدرجتي حرية، والقيمة $d^2>9.21$ لا تحدث إلا في $1\%$ من الإطارات. السياسة \en{none}: لا رفض أبداً. السياسة \en{hybrid}:
يُرفض القياس المفاجئ، وبعد $3$ رفضات متتالية يُعاد ضبط المرشح على القياس مع $P=I$.

\subsection{المتغيرات}
\begin{itemize}
\item \en{cvkf}: المرشح نفسه مع $c=0$ وبلا جدران ولا ذراع (مرشح كالمان بسرعة ثابتة).
\item \en{ekf19}: احتكاك $c=0.85$ و$E=0.725$ (منتصفا مدى السحب).
\item \en{ekf19b}: \en{ekf19} مع عدم يقين إضافي عند الارتداد: إذا غيّرت مركبة السرعة إشارتها يُضاف $(\sigma_E|v|)^2$ إلى تباينها،
حيث $\sigma_E=0.13$ هو الانحراف المعياري للتوزيع $U(0.5,0.95)$: $0.45/\sqrt{12}=0.13$.
\item \en{ekf19b\_fit} (جديد، \en{B3}): \en{ekf19b} مع $c,E$ ملاءمَين على التدريب (الخطوة $1$).
\item \en{oracle} (العرّاف): المرشح نفسه مع $\mu,E$ \emph{الحقيقيين} لكل رمية؛ حدّ أعلى لأي مرشح بثوابت ثابتة.
\item \en{+stale} (\en{StaleGuard}): مُغلِّف يعتبر القياس الصالح المطابق تماماً للسابق مفقوداً.
\end{itemize}

%=====================================================================================================
\section{الخلفية الرابعة: كيف تُقاس الأخطاء}

\subsection{متى وماذا (الملف \code{estimators/metrics.py})}
يُقيَّم المرشح في لحظة \textbf{الوصول} $t_a=t_{\text{capture}}+\text{delay}$، وهي اللحظة التي يملك فيها الروبوت المعلومة فعلاً:
\begin{itemize}
\item \en{now\_pos}: خطأ الموقع في $t_a$؛ و\en{now\_vel}: خطأ السرعة (م/ث).
\item \en{h0.2} و\en{h0.5} و\en{h1.0}: خطأ موقع التنبؤ بعد $0.2$ و$0.5$ و$1.0$ ث من $t_a$ مقارنة بالموقع الحقيقي حينها.
\en{h0.5} هو الرقم الرئيسي، لأن الإمساك بالكرة يحتاج تنبؤاً بهذا الأفق تقريباً.
\end{itemize}
\textbf{أي إطارات تُقيَّم:} كل إطار سادس (عشرة في الثانية)، بعد $0.1$ ث من الإحماء فقط، وفقط ما دامت سرعة الكرة الحقيقية
$\ge0.05$ م/ث، والتنبؤ فقط إذا كان زمنه داخل التسجيل. الحقيقة هي مسار المحاكي بتردد $500$ هرتز مع استيفاء خطي.

\subsection{طريقة جديدة \#4: الإحصاء على مستوى الرمية (بند \en{A3})}
\textbf{لماذا لا نجمع كل الإطارات معاً؟}
\begin{example}[تجميع الإطارات مضلِّل]
الرمية بسرعة $5$ م/ث قد تُقيَّم على نحو $90$ إطاراً، والرمية بسرعة $0.1$ م/ث على $1$ إلى $3$ إطارات فقط. إذا جمعنا كل الإطارات
في رقم واحد، فالرميات السريعة هي التي تقرر النتيجة: $69\%$ من كل الإطارات المقيَّمة تأتي من رميات سرعتها $\ge2.5$ م/ث.
ثم إن إطارات الرمية الواحدة مترابطة بقوة (الكرة نفسها و$\mu$ نفسها)، فهي ليست $90$ دليلاً مستقلاً بل أقرب إلى دليل واحد.
\end{example}

\textbf{مثال رقمي صغير.} رميتان: الأولى سريعة، $90$ إطاراً كلها بخطأ $50$ مم؛ والثانية بطيئة، $2$ إطاران بخطأ $200$ مم.
\begin{itemize}
\item الطريقة القديمة (تجميع الإطارات): $\sqrt{(90\times50^2+2\times200^2)/92}=\sqrt{(225000+80000)/92}=57.6$ مم، فالرمية البطيئة
السيئة شبه غير مرئية.
\item الطريقة الجديدة (لكل رمية): $\mathrm{RMS}_1=50$ و$\mathrm{RMS}_2=200$، والمتوسط على الرميات $125$ مم، فكل رمية لها صوت واحد.
\end{itemize}

\paragraph{\en{RMS} لكل رمية.} للرمية $i$ ذات الأخطاء $e_{i1},\dots,e_{in_i}$:
\[ \mathrm{RMS}_i=\sqrt{\tfrac{1}{n_i}\textstyle\sum_k e_{ik}^2}. \]
تعرض الجداول \textbf{متوسط الرميات} $\bar R=\frac1N\sum_i\mathrm{RMS}_i$ و\textbf{وسيط الرميات}، ومعهما للاستمرارية الرقم
القديم المجمَّع. ولماذا \en{RMS} لا المتوسط العادي؟ لأن \en{RMS} يعاقب الأخطاء الكبيرة أكثر، والخطأ الكبير هو الذي يُفشل الإمساك.
ولماذا نعرض الوسيط أيضاً؟ لأنه لا يتأثر برمية واحدة شاذة، فهو يصف الحالة «النموذجية».

\subsection{طريقة جديدة \#5: فاصل الثقة بطريقة \en{bootstrap}}
\textbf{السؤال:} لدينا $N=75$ رمية فقط. لو أعدنا التجربة برميات أخرى، كم سيتغير $\bar R$؟

\textbf{الإجراء:}
\begin{enumerate}
\item اسحب $N$ رمية \emph{مع الإرجاع} من الرميات الـ$N$ التي لدينا (بعضها يُسحب مرتين وبعضها لا يُسحب).
\item احسب $\bar R$ لهذه المجموعة الجديدة.
\item كرر $2000$ مرة. النقطتان $2.5\%$ و$97.5\%$ من القيم الـ$2000$ تشكلان فاصل الثقة $95\%$.
\end{enumerate}

\textbf{الفكرة:} رمياتنا الـ$75$ عيِّنة من كل الرميات الممكنة. أفضل ما نعرفه عن «كل الرميات الممكنة» هو عيّنتنا نفسها، فإعادة
السحب منها تحاكي سحب عينات جديدة من العالم. الفاصل العريض يعني «لا نستطيع الجزم بدقة».

\begin{example}[بوتستراب على $4$ رميات]
لتكن أخطاء الرميات $(40,\,50,\,60,\,250)$ مم، والمتوسط $100$. بعض السحوبات الممكنة:
$(40,40,50,60)\to47.5$، و$(50,250,250,60)\to152.5$، و$(40,50,60,60)\to52.5$. القيم تتراوح كثيراً لأن الرمية الشاذة ($250$) قد
تُسحب مرتين أو لا تُسحب أبداً. الفاصل الناتج عريض، وهذا صادق: أربع رميات لا تكفي.
\end{example}

\subsection{طريقة جديدة \#6: المقارنة المقترنة واختبار ويلكوكسون}
\textbf{الاقتران.} تُشغَّل الطريقتان \en{A} و\en{B} على الرميات \emph{نفسها} وبالضوضاء \emph{نفسها}. لكل رمية نأخذ الفرق
$d_i=\mathrm{RMS}_i^A-\mathrm{RMS}_i^B$ (السالب يعني أن \en{A} أفضل).

\textbf{لماذا يقوي الاقتران المقارنة؟ (رياضياً)}
\[ \mathrm{Var}(A-B)=\mathrm{Var}(A)+\mathrm{Var}(B)-2\,\mathrm{Cov}(A,B). \]
الرميات تختلف كثيراً فيما بينها (سريعة/بطيئة، ارتدادات كثيرة/قليلة)، فالطريقتان كلتاهما تخطئان أكثر على الرميات الصعبة،
أي أن $\mathrm{Cov}(A,B)$ كبير وموجب. الطرح يزيل هذا التباين المشترك، فيبقى الفرق الحقيقي بين الطريقتين، وتصبح الفروق الصغيرة
الحقيقية مرئية. لو قارنّا متوسطَي مجموعتين مستقلتين لضاع الفرق في ضجيج اختلاف الرميات.

\textbf{ما يُعرض:}
\begin{itemize}
\item \textbf{متوسط الفرق} مع فاصل \en{bootstrap}: إذا لم يحتوِ الفاصل على الصفر فالفرق حقيقي.
\item \textbf{اختبار ويلكوكسون للرتب الموقَّعة} (\en{Wilcoxon signed-rank})، وفقراته:
\begin{enumerate}
\item رتِّب القيم $|d_i|$ من الأصغر إلى الأكبر.
\item اجمع رتب الفروق الموجبة، وسمِّ المجموع $W^+$.
\item لو كانت الطريقتان متساويتين، لكانت إشارة كل فرق مثل رمي عملة، ولكان $W^+$ قريباً من نصف مجموع الرتب.
\item القيمة الاحتمالية \en{p} هي احتمال الحصول على $W^+$ متطرف بهذا القدر لو كانتا متساويتين.
\end{enumerate}
القيمة الصغيرة (مثلاً $<0.01$) تعني أن الفرق يصعب جداً أن يكون حظاً. ولأن الاختبار يستخدم الرتب لا القيم، لا تستطيع رمية
واحدة مجنونة أن تسيطر عليه.
\item \textbf{نسبة الفوز}: نسبة الرميات التي كانت فيها \en{A} أفضل.
\end{itemize}

\begin{example}[ويلكوكسون على $5$ رميات]
لتكن $d=(-12,-3,+1,-8,-20)$ مم.
\begin{itemize}
\item رتب $|d|$ هي $(4,2,1,3,5)$، والفرق الموجب الوحيد هو $+1$ ورتبته $1$، فإن $W^+=1$.
\item لو كانت الطريقتان متساويتين لكانت كل إشارة رمية عملة، فعدد أنماط الإشارات $2^5=32$.
\item نمطان فقط يعطيان $W^+\le1$: لا فروق موجبة، أو الأصغر وحده موجب. فالاحتمال من جهة واحدة $2/32$، ومن الجهتين $4/32=0.125$.
\end{itemize}
خمس رميات لا تكفي للجزم، ولهذا نحتاج $75$.
\end{example}

\begin{warn}[المقارنات الكثيرة]
الخطوة $4$ تجري عشرات المقارنات. عند العتبة $p<0.05$ يظهر نحو واحد من كل عشرين نتيجة «دالّة» بالصدفة حتى لو لم يكن هناك
فرق. لذلك نفسّر في الأطروحة فقط المقارنات المخطَّطة مسبقاً (المذكورة في الخطوة $4$)، ونعرض الفواصل لا القيم \en{p} وحدها،
ونستخدم عتبة صارمة مثل $p<0.01$.
\end{warn}

%=====================================================================================================
\section{الخطوة 1: ملاءمة $c$ و$E$ على التدريب}
\noindent\textbf{الملف:} \code{step\_0k\_fit\_constants\_train.py}\par\medskip
\begin{why}[الغرض (بند \en{B3}): طريقة جديدة \#7، المقارنة العادلة]
كان \en{ekf19b} يستخدم $c=0.85$ و$E=0.725$، وهما منتصفا المدى الذي \emph{سُحبت} منه الرميات؛ أي معرفة بمولِّد البيانات لا يملكها
نظام حقيقي. المرشحات المتعلِّمة في المرحلة الأولى ستلائم أرقامها على رميات التدريب، فالمقارنة العادلة يجب أن تلائم ثوابتها
على الرميات نفسها وبالمعيار نفسه. وإلا فقد يأتي أي تحسّن للطريقة الجديدة من «ضبط أكثر» لا من الفكرة نفسها.
\end{why}
\textbf{الإجراء:}
\begin{enumerate}
\item لكل زوج $(c,E)$ من الشبكة
\[ c\in\{0.45,0.55,\dots,1.05\}\ (7\text{ قيم})\ \times\ E\in\{0.625,0.675,\dots,0.825\}\ (5\text{ قيم}) \]
شغِّل \en{ekf19b-none+stale} على كل رميات التدريب الـ$350$.
\item الضوضاء $(\sigma,q)$ مأخوذة من ضبط التحقق في المرحلة صفر.
\item النتيجة لكل زوج هي متوسط \en{h0.5} على الرميات، والأقل يفوز.
\item يُجرى هذا في الحالة \en{ideal} لأن الثوابت الفيزيائية لا تعتمد على الكاميرا، ويُكرَّر في \en{all\_real} للتحقق فقط.
\end{enumerate}
\textbf{لماذا قد يختلف أفضل $c$ عن $0.85$؟} خطأ التنبؤ ليس متماثلاً حول $c$: المرشح ذو $c$ الكبير جداً يتنبأ بتوقف مبكر،
فيكبر خطؤه مع بقية المسافة. والرميات ذات $\mu$ الكبير تتوقف أسرع فتُقيَّم على إطارات أقل، لكن الإحصاء لكل رمية يعطي كل رمية
وزناً متساوياً. لهذا قد يكون أفضل ثابت أدنى أو أعلى من المتوسط، واكتشاف ذلك هو هدف الخطوة.

\textbf{النزول الإحداثي لمرة واحدة (\en{coordinate descent}).} نلائم $(c,E)$ مع تثبيت $(\sigma,q)$، ثم تعيد الخطوة $2$ ضبط
$(\sigma,q)$ مع $(c,E)$ الملاءمَين. الملاءمة المشتركة الكاملة كانت ستكرر هذا التناوب حتى يستقر. تمريرة واحدة شائعة، ونصرّح
بذلك. (النتيجة الفعلية في القسم الأخير تبيّن أن هذا كان مهماً.)

\textbf{المخرجات:} \code{outputs/fit\_constants\_train.json} والسجل \code{fit\_constants\_train.log}.

%=====================================================================================================
\section{الخطوة 2: ضبط الضوضاء على التحقق}
\noindent\textbf{الملف:} \code{step\_0i\_tune\_ekf\_noise.py}\par\medskip
\begin{why}[الغرض]
$\sigma$ (مقدار الثقة بالكاميرا) و$q$ (مقدار الثقة بنموذج الحركة) يغيّران النتائج كثيراً، وأفضل قيمهما تعتمد على حالة الحساس.
الضبط على رميات \emph{التحقق} يُبقي التدريب والاختبار نظيفين. الجديد هنا:
\begin{itemize}
\item كل عائلة تُضبط وحدها، فتحصل كل طريقة على جهد الضبط نفسه.
\item المعيار متوسط الرميات لا الإطارات المجمَّعة.
\item الحالات الجديدة مشمولة.
\end{itemize}
\end{why}
\begin{center}\small
\begin{tabularx}{\linewidth}{@{}l>{\raggedleft\arraybackslash}Xll@{}}\toprule
العائلة & تُضبط باستخدام & شبكة $\sigma$ (مم) & شبكة $q$\\\midrule
\en{ekf} & \en{ekf19b-none+stale} ($c,E$ المنتصفان) & $2,3,5,8,12,20,30$ & $0.03,0.1,0.5,2,5,20$\\
\en{ekf\_fit} & \en{ekf19b\_fit-none+stale} ($c,E$ من الخطوة $1$) & كما سبق & كما سبق\\
\en{oracle} & \en{ekf19b\_oracle-none+stale} ($\mu,E$ الحقيقيان) & كما سبق & كما سبق\\
\en{cvkf} & \en{cvkf} & $1,2,3,5,8,12$ & $2,5,20,50,200,1000$\\\bottomrule
\end{tabularx}\end{center}
\begin{itemize}
\item \textbf{حجم العمل:} $162$ إعداداً $\times75$ رمية تحقق $\times11$ حالة، أي نحو $134\,000$ تشغيل مرشح.
\item \textbf{المعيار:} متوسط \en{h0.5} على رميات التحقق، والأقل يفوز لكل عائلة ولكل حالة.
\item \textbf{لماذا شبكة \en{cvkf} خاصة وأوسع؟} لأنه بلا احتكاك ولا جدران، فيحتاج $q$ كبيراً جداً ليتبع التباطؤ والارتدادات.
في المرحلة صفر كان أفضل $q$ له على حافة الشبكة في $5$ من $8$ حالات.
\item \textbf{لماذا يُضبط العرّاف أيضاً؟} الحدّ الأعلى يجب أن يكون \emph{أفضل} ما تسمح به الثوابت الحقيقية. في المرحلة صفر أُعطي
إعدادات \en{ekf19b} بدلاً من إعداداته الخاصة.
\end{itemize}

\begin{idea}[مفهوم مهم: «حافة الشبكة»]
البحث الشبكي يجرّب قيماً محددة فقط. إذا كانت أفضل قيمة هي الأصغر أو الأكبر في الشبكة، فقد تكون القيمة الأفضل فعلاً خارجها،
ونحن لم نجربها. لذلك يطبع الكود علامة \en{ON GRID EDGE}، والقاعدة البحثية: النتيجة على الحافة ليست نهائية، ويجب توسيع الشبكة
وإعادة الضبط.
\end{idea}

\textbf{المخرجات:} \code{outputs/ekf\_tuning\_val.json} والسجل الذي يطبع كل الشبكات. نسخة المرحلة صفر محفوظة باسم
\code{ekf\_tuning\_val\_v2\_phase0.json}.

%=====================================================================================================
\section{الخطوة 3: المقارنات الأساسية}
\noindent\textbf{الملف:} \code{step\_0d\_baselines\_with\_arm.py}\par\medskip
\begin{itemize}
\item \textbf{التشغيل:} $13$ مقدِّراً على $75$ رمية تحقق و$75$ رمية اختبار تطويري، في كل الحالات الـ$11$، وكلٌّ بضوضاء عائلته
المضبوطة (و\en{ekf19b\_fit} بثوابت الخطوة $1$ أيضاً).
\item \textbf{بذرة الضوضاء} للرمية $i$ هي $1000+i$ في كل مكان، فيرى كل مقدِّر القياسات المشوَّشة نفسها تماماً.
\item \textbf{الجداول:} لكل حالة ولكل من التحقق والاختبار، وفي كل خلية: متوسط الرميات [فاصل $95\%$]، ثم وسيط الرميات، ثم
\en{RMS} المجمَّع القديم.
\item \textbf{المخرجات:}
\begin{itemize}
\item \code{baselines\_v3\_val\_test.json}: الملخصات.
\item \code{baselines\_v3\_errors.pkl.gz}: \emph{كل} خطأ لكل إطار ولكل مقدِّر ورمية وحالة، فيمكن حساب أي إحصاء جديد لاحقاً
دون إعادة التشغيل.
\end{itemize}
\item \textbf{لماذا التحقق والاختبار التطويري فقط، لا الـ$500$ كلها؟} أرقام «الكل» في المرحلة صفر كانت تخلط رميات التدريب،
والخطوة $1$ تلائم عليها الآن. العرض على بيانات استُخدمت في الملاءمة يجامل الطريقة الملاءَمة.
\end{itemize}

%=====================================================================================================
\section{الخطوة 4: الاختبارات المقترنة}
\noindent\textbf{الملف:} \code{step\_0j\_paired\_tests.py}\par\medskip
تستخدم الأخطاء المحفوظة من الخطوة $3$ دون أي تشغيل جديد. لكل من الاختبار التطويري والتحقق، ولكل من \en{h0.5} و\en{now\_pos}:
متوسط الفرق مع فاصله، وقيمة \en{p} لويلكوكسون، ونسبة الفوز.
\begin{center}\small
\begin{tabularx}{\linewidth}{@{}>{\raggedleft\arraybackslash}p{6.2cm}>{\raggedleft\arraybackslash}X@{}}\toprule
المقارنة & السؤال\\\midrule
\en{ekf19b\_fit} مقابل \en{ekf19b} & هل الثوابت الملاءَمة على التدريب أفضل من منتصفَي المدى؟\\
\en{oracle} مقابل \en{ekf19b} & ما حجم الفجوة التي تحاول المرحلة الأولى سدّها؟\\
\en{ekf19b-hybrid} مقابل \en{ekf19b+stale} & هل البوابة تفيد أم تضر؟\\
\en{ekf19b+stale} مقابل \en{cvkf} & كم يضيف النموذج الفيزيائي؟\\\midrule
\multicolumn{2}{@{}r}{بين الحالات (لـ \en{ekf19b+stale} وللعرّاف):}\\
\en{plane} مقابل \en{noise\_real} & \en{B1}: هل قياس المستوى أفضل من عمق شبيه بـ \en{D435}؟\\
\en{all\_plane} مقابل \en{all\_real} & \en{B1} في الحالة الواقعية الكاملة\\
\en{noise} مقابل \en{ideal}؛ \en{noise\_real} مقابل \en{noise} & كم تكلّف كل درجة من ضوضاء الكاميرا؟\\
\en{dropout\_silent\_noisy} مقابل \en{noise} & \en{B4}: كم تكلّف الإطارات القديمة غير القابلة للكشف؟\\\bottomrule
\end{tabularx}\end{center}

%=====================================================================================================
\section{الخطوة 5: $\mu$ مقابل $E$ على التدريب والتحقق}
\noindent\textbf{الملف:} \code{step\_1a\_mu\_vs\_e\_gap.py}\par\medskip
\begin{why}[الغرض (بند \en{A2})]
قرار «أدخل $\mu$ في الحالة أولاً و$E$ لاحقاً» اتُّخذ من فحص على الرميات الـ$500$ كلها \emph{بما فيها} رميات الاختبار. يُعاد الفحص
على التدريب والتحقق فقط ($425$ رمية) وبإحصاء الرميات. إذا صمد الاستنتاج صار القرار نظيفاً، وإذا تغيّر وجب إعادة النظر في
المرحلة الأولى.
\end{why}
أربعة مرشحات في الحالتين \en{ideal} و\en{all\_real}: الاسمي ($c=0.85$، $E=0.725$)؛ $\mu$ الحقيقي فقط؛ $E$ الحقيقي فقط؛ وكلاهما
حقيقيان (أي العرّاف). لكل منها: متوسط الرميات وفاصله، والوسيط، والفرق المقترن مع الاسمي.

%=====================================================================================================
\section{الخطوة 6: فحص الفيزياء الواقعية}
\noindent\textbf{الملف:} \code{step\_0l\_check\_mismatch\_physics.py}\par\medskip

\subsection{طريقة جديدة \#8: لماذا فيزياء ثانية أصلاً؟ (بند \en{A1})}
\textbf{الجريمة العكسية} (\en{inverse crime}) مصطلح من علوم المسائل العكسية: أن تختبر طريقتك على بيانات ولّدتها بالنموذج نفسه
الذي تفترضه الطريقة. النتيجة تبدو رائعة لكنها لا تقول شيئاً عن الواقع. عندنا: خطأ العرّاف النموذجي على $0.5$ ث نحو
$10$ مم فقط، لأن نموذجه هو الحقيقة حرفياً.

والأخطر: المرحلة الأولى ستضع $\mu$ ثم $E$ في حالة المرشح. وبما أن الحقيقة في المحاكي هي بالضبط «$\mu$ ثابت و$E$ ثابت لكل رمية»،
فهذه الطريقة مضمونة النجاح مسبقاً. والمرحلة الثانية ستقارن المرشحات المتعلِّمة (الصندوق الأسود) بالرمادية في عالم \emph{هو}
النموذج الرمادي نفسه، فالنتيجة محسومة بالمولِّد لا بالطرق.

\textbf{الحل:} محاكٍ ثانٍ (\code{MismatchArena}) يضيف ظواهر معروفة في فيزياء الكرات المتدحرجة لا يحتويها \emph{أي} مقدِّر.
الاستنتاج يُعتمد فقط إذا صمد على المجموعتين.

\subsection{الفيزياء الواقعية بالرياضيات}
للكرة الآن دوران $\bm\omega$ (ثلاثة أرقام) يُتتبَّع بالبايثون. التصيير لا يتغير، فالكرة الخضراء الموحدة اللون تبدو نفسها وهي تدور.
الكرة كرة صلبة عزم قصورها $I=\frac25mr^2$.

\paragraph{(أ) الانزلاق ثم التدحرج.} تلامس الأرض الكرة في النقطة $-r\hat z$ تحت المركز، وسرعة هذه النقطة هي \emph{سرعة الانزلاق}:
\[ \uu=\vv+\bm\omega\times(-r\hat z)=(v_x-r\omega_y,\ v_y+r\omega_x). \]
\begin{itemize}
\item $\uu=0$ يعني \emph{تدحرجاً} (لا انزلاق عند التلامس).
\item ما دامت $\uu\ne0$ فالكرة \emph{تنزلق}، ويؤثر احتكاك الأرض $\mu_sg$ عكس $\uu$:
\[ \dot\vv=-\mu_sg\,\hat\uu,\qquad \dot{\bm\omega}=\frac{\bm r_c\times\bm F}{I}=\frac{5\mu_sg}{2r}\,(\hat z\times\hat\uu). \]
\item نعوّضهما في $\dot\uu=\dot\vv-r\,\dot{\bm\omega}\times\hat z$، ومع $(\hat z\times\hat\uu)\times\hat z=\hat\uu$ لأي $\hat\uu$ أفقي:
\[ \dot\uu=-\mu_sg\,\hat\uu-\tfrac52\mu_sg\,\hat\uu=-\tfrac72\,\mu_sg\,\hat\uu. \]
\item الانزلاق يتناقص بمعدل ثابت $\frac72\mu_sg$، ويبدأ التدحرج بعد $t_r=\dfrac{2|\uu_0|}{7\mu_sg}$.
\item لكرة أُطلقت بلا دوران ($\uu_0=\vv_0$) تكون عندها $|\vv|=v_0-\mu_sg\,t_r=v_0-\tfrac27v_0=\tfrac57v_0$.
\end{itemize}

\textbf{برهان ثانٍ بالعزم الزاوي.} الاحتكاك يؤثر في نقطة التلامس، فلا يغيّر العزم الزاوي حولها:
\[ mrv_0=mrv_f+I\frac{v_f}{r}\ \Rightarrow\ v_f=\frac{v_0}{1+2/5}=\frac57v_0. \]

\textbf{المقدِّر لا يعرف هذا}: يفترض تباطؤاً ثابتاً واحداً، بينما الحقيقة تتباطأ بقوة ($\mu_sg\approx1.5$ إلى $4$ م/ث$^2$) لفترة قصيرة
ثم بضعف. بلغة المرحلة الأولى: «$\mu$ الفعّال» يتغير أثناء الرمية، وسنرى إن كان \en{ekf-$\mu$} يتعلم ذلك عبر $q_\mu$.

\paragraph{(ب) مقاومة التدحرج.} أثناء التدحرج يكون التباطؤ $a_{\text{roll}}=a_0(1+k|\vv|)$ عكس $\vv$، ويتبع الدوران شرط التدحرج.
القيم المنشورة للسجاد: معامل مقاومة تدحرج نحو $0.013$ إلى $0.047$، أي $a_0\approx0.1$ إلى $0.5$ م/ث$^2$، واخترتَ المدى $0.1$ إلى $0.8$.

\paragraph{(ج) دوران الإطلاق.} الانزلاق الابتدائي $\uu_0=s\,\vv_0$ مع $s\sim U(0,1)$: القيمة $s=0$ تعني تدحرجاً من البداية (كرة تُترك
من منحدر)، و$s=1$ انزلاقاً خالصاً (كرة مضروبة).

\paragraph{(د) ميل الأرض.} ميل $\alpha$ في اتجاه $\phi$ يضيف $\bm g_t=g\sin\alpha\,(\cos\phi,\sin\phi)$. للكرة \emph{المتدحرجة} لا يسرّع
المركز إلا $\frac57$ منه:
\[ ma=mg_t-F,\quad I\frac{a}{r}=Fr\ \Rightarrow\ a\left(1+\frac{I}{mr^2}\right)=g_t\ \Rightarrow\ a=\frac{g_t}{1.4}=\frac57g_t. \]
(جزء من القوة يذهب لتسريع الدوران.) تبقى الكرة ساكنة إذا $\frac57g_t<a_0$. ومع $\alpha\le0.15^\circ$:
$\frac57g\sin\alpha\le0.018<0.1$، فالكرة الساكنة لا تبدأ التدحرج وحدها أبداً.

\paragraph{(هـ) الاصطدام بالجدران والذراع.} العمودي $\bm n$ يشير إلى داخل الساحة (أو خارج صندوق الذراع).
\begin{itemize}
\item \textbf{العمودي}: $v_n'=-e(v_n)\,v_n$ مع معامل ارتداد يتناقص مع سرعة الاصطدام:
$e=\mathrm{clip}(e_0-\beta|v_n|,\,0.3,\,0.98)$، و$e_0$ منفصل للجدران ($U(0.55,0.95)$) وللذراع ($U(0.40,0.90)$).
\item \textbf{الدفع العمودي} لوحدة الكتلة: $J_n=(1+e)|v_n|$.
\item \textbf{الاحتكاك عند التلامس}: التلامس على ارتفاع مركز الكرة.
\begin{itemize}
\item انزلاق نقطة التلامس في مستوى الجدار له جزء أفقي (على طول الجدار) وجزء رأسي (من دوران التدحرج).
\item دفع كولوم (\en{Coulomb}) يعاكسه بمقدار $\min(\mu_wJ_n,\ \text{الدفع الذي يوقف الانزلاق})$، مع $\mu_w\sim U(0.1,0.3)$
(قاس \en{Mathavan} وزملاؤه $2010$ القيمة $0.14$).
\item الجزء الأفقي يغيّر السرعة المماسية، والرأسي يغيّر الدوران فقط، لأن الأرض تمسك الكرة.
\end{itemize}
\item \textbf{التوقيت الدقيق}: تُعاد الكرة إلى نقطة الاصطدام (زمن $\tau=\text{التجاوز}/|v_n|$ قبل نهاية الخطوة)، وتُغيَّر سرعتها،
ثم تُحرَّك للأمام بالسرعة الجديدة؛ أي عكس لا قصّ (إصلاح \en{C4} في هذه البيانات).
\end{itemize}

\begin{example}[أكبر فرق: كرة متدحرجة تصطدم بالجدار وجهاً لوجه]
\begin{itemize}
\item قبل الاصطدام تتدحرج الكرة بسرعة $v$ ودوران $\omega=v/r$ (دوران أمامي).
\item بعد الاصطدام تصبح السرعة $-ev$، لكن الدوران ما زال «أمامياً»، فتنزلق الكرة ويبطئ الانزلاق عودتها.
\item احتكاك الجدار يزيل أولاً نسبة $f=\frac52\mu_w(1+e)$ من الدوران (إن لم يتوقف الانزلاق كلياً).
\item بحفظ العزم الزاوي حول نقطة تلامس الأرض تتدحرج الكرة أخيراً عائدةً بسرعة
\[ v_{\text{final}}=\frac{-ev+\frac25(1-f)v}{1+\frac25}=-\frac{5e-2(1-f)}{7}\,v. \]
\end{itemize}
مع $e=0.8$: إذا $f=0$ نحصل على $0.29\,v$، ومع $\mu_w=0.14$ ($f=0.63$) نحصل على $0.47\,v$. بينما تفترض المقدِّرات $0.8\,v$.
هذا خطأ نموذج كبير وواقعي، وهو بالضبط ما صُممت هذه البيانات لكشفه.
\end{example}

\paragraph{الاقتران بين المجموعتين.} الرميات الواقعية تحتفظ بنقطة البداية والسرعة \emph{نفسيهما} للرميات المطابقة (البذرة $42$)،
والأرقام الفيزيائية الإضافية من تيار عشوائي منفصل (البذرة $43$). فالمجموعتان لا تختلفان إلا في الفيزياء، رمية برمية
(تحققنا: الفرق في $\pp_0$ و$\vv_0$ صفر).

\paragraph{العرّاف على البيانات الواقعية.} لا يوجد هنا $\mu$ أو $E$ «حقيقي» واحد، فالعرّاف القديم لا معنى له. قرارك: العرّاف
بالنظر إلى الوراء (\en{hindsight}) هو الثابتان $(c,E)$ اللذان يلائمان أفضل ملاءمة المسار الحقيقي لكل رمية، أي أفضل ما يستطيعه
أي مرشح بثوابت ثابتة.

\subsection{ما تفحصه الخطوة 6}
\textbf{$12$ فحصاً مقابل معادلات مغلقة}: التدحرج عند $\frac57v_0$ وزمنه؛ تباطؤ التدحرج $a_0(1+kv)$؛ الميل $\frac57g\sin\alpha$؛
ارتداد الجدار $e_0$ و$e_0-\beta|v_n|$ مع ثبات المماسي؛ حفظ الدوران بلا احتكاك جدار، وإزالة النسبة $\frac52\mu_w(1+e)$ معه؛
وسرعة التدحرج النهائية $-(5e-2(1-f))/7$.

\textbf{على الرميات الـ$500$} (الحقيقة فقط، بلا تصيير): لا اصطدام يزيد الطاقة الحركية (انتقالية ودورانية)؛ لا تنتهي خطوة والكرة
داخل الذراع أو خلف جدار؛ وإحصاءات مقارنة بين المطابق والواقعي.

%=====================================================================================================
\section{الخطوة 7: تصيير البيانات الواقعية (أُجِّلت بطلبك)}
\noindent\textbf{الأمر:} \code{step\_0c\_generate\_500\_with\_arm.py --physics mismatch}\par\medskip
\begin{itemize}
\item \textbf{العمل:} يحاكي كل رمية من الرميات الـ$500$ المقترنة لمدة $10$ ث بالفيزياء الواقعية، ويصيّر الكاميرا بتردد $60$ هرتز (ألوان
وعمق، $480\times480$)، ويخزن الكشوفات المثالية كما في المجموعة المطابقة.
\item \textbf{الزمن:} نحو $35$ ثانية حساب لكل رمية، أي نحو $90$ دقيقة بثلاث عمليات.
\item \textbf{المخزَّن:} في كل ملف: الحقيقة بتردد $500$ هرتز، والكشوفات، وسجل الارتدادات، وكل الأرقام الفيزيائية تحت المفاتيح
\code{phys\_*}.
\item \textbf{$\mu$ و$e$ يُخزَّنان كقيمة \en{NaN}} عمداً، كي لا يُشغَّل العرّاف القديم بالخطأ.
\item \textbf{قابل للاستئناف:} الملفات الجاهزة تُتخطى، وكل ملف يُكتب باسم مؤقت أولاً، فالانقطاع لا يترك ملفاً نصف مكتوب.
\end{itemize}
\textbf{للتشغيل لاحقاً:} احذف العلامة ثم شغّل الأمر:
\begin{center}\code{rm outputs/.overnight\_done/render\_mismatch\_data}\\[2pt]
\code{bash experiments/run\_overnight.sh part1}\end{center}

%=====================================================================================================
\section{نتائج التشغيل الفعلي (ليلة \LR{2026-11-01})}
انتهى الجزء الأول الساعة $01{:}29$ بلا أخطاء (التصيير أُجِّل). التفاصيل الكاملة في الملف:
\begin{center}\code{docs/part1\_results\_2026-11-02.md}\end{center}
الأرقام على الاختبار التطويري: متوسط الرميات لخطأ \en{h0.5}، وبين القوسين فاصل $95\%$.

\begin{enumerate}
\item \textbf{ملاءمة $c,E$ لا تفيد}: أفضل ملاءمة $(0.75,\,0.725)$ تعطي $51.2$ مم مقابل $51.4$ للمنتصفين، والاختبار المقترن لا يجد فرقاً
في أي حالة ($p\ge0.3$). أي أن المنتصفين كانا شبه مثاليين، وبند \en{B3} صار مقيساً لا مفترضاً: أثره ضئيل. (وبعد إعادة ضبط
الضوضاء صار $0.85$ هو الأفضل؛ هذا أثر النزول الإحداثي لمرة واحدة.)
\item \textbf{المقارنة الأساسية مقابل العرّاف}: \en{ekf19b} في الحالة المثالية $54\,[47,62]$ مم والعرّاف $7$ مم؛ وفي
\en{all\_real} $76\,[67,84]$ مقابل $20$. الفجوة المقترنة بين $-45$ و$-58$ مم، والعرّاف أفضل في $95$ إلى $100\%$ من الرميات.
\item \textbf{المستوى أفضل من العمق (\en{B1})}: أفضل بمقدار $9.8\,[7.6,11.7]$ مم، وفي $95\%$ من الرميات، و$p=4\times10^{-11}$،
وفي الحالة الكاملة $10.4$ مم. اعتراض المحكِّم كان صحيحاً.
\item \textbf{الإطارات القديمة غير القابلة للكشف (\en{B4})}: تضيف $36.5\,[22.9,52.7]$ مم، وحتى العرّاف يخسر $34$ مم.
\item \textbf{البوابة تضر}: من $8$ إلى $11$ مم في الحالات العادية، و$+62$ مم مع الإطارات المجمَّدة.
\item \textbf{«$\mu$ أولاً» مؤكَّد بلا بيانات اختبار}: معرفة $\mu$ الحقيقي توفر $21$ مم (أفضل في $87\%$ من الرميات)، ومعرفة
$E$ فقط $11$ مم ($60\%$).
\item \textbf{الفيزياء الواقعية سليمة}: $12/12$ فحصاً، صفر اصطدامات تزيد الطاقة، والصعوبة مشابهة (ارتدادات لكل رمية $1.6$ مقابل
$2.0$، وزمن الحركة الوسيط $2.6$ مقابل $2.3$ ث).
\end{enumerate}

\begin{warn}[مشكلة يجب إصلاحها قبل اعتماد الأرقام: حافة الشبكة]
\begin{itemize}
\item \en{ekf19b}: أفضل $\sigma=2$ مم (أصغر قيمة) في $6$ حالات.
\item العرّاف: $\sigma=30$ مم (أكبر قيمة) و$q=0.03$ (أصغر قيمة) في كل الحالات. لأن نموذجه كامل، فهو يريد تجاهل الكاميرا تقريباً،
وهذه علامة أخرى على الجريمة العكسية.
\item \en{cvkf}: $q=1000$ (أكبر قيمة).
\end{itemize}
الإصلاح المقترح: توسيع الشبكات وإعادة الضبط (نحو $2.5$ ساعة)، ثم إعادة المقارنات والاختبارات (نحو $20$ دقيقة). بعدها تصبح أرقام
الجزء الأول نهائية.
\end{warn}

%=====================================================================================================
\section{التشغيل والإيقاف والاستئناف}
\begin{codebox}[الأوامر (من مجلد \en{ball\_estimation})]
\begin{itemize}
\item لتشغيل الجزء الأول وحده:\\ \code{bash experiments/run\_overnight.sh part1}
\item لتشغيل الجزأين معاً:\\ \code{bash experiments/run\_overnight.sh}
\item التقدّم يُكتب في \code{outputs/overnight\_progress.log}.
\item كل خطوة منتهية تترك علامة في \code{outputs/.overnight\_done/}، وإعادة الأمر نفسه تتخطى الخطوات المنتهية. لإعادة كل شيء
احذف هذا المجلد.
\end{itemize}
\end{codebox}
\begin{warn}[حدود الجهاز]
ثلاث عمليات كحد أقصى ($4$ أنوية، $7$ جيجابايت ذاكرة)، فعدد أكبر تسبب على الأرجح في إطفاء سابق. يجب أن يبقى الحاسوب مشغَّلاً
وموصولاً بالشاحن. إذا انطفأ لا يضيع شيء منتهٍ.
\end{warn}

%=====================================================================================================
\section{مسرد المصطلحات}
\begin{center}\small
\begin{tabularx}{\linewidth}{@{}l>{\raggedleft\arraybackslash}X@{}}\toprule
المصطلح & المعنى\\\midrule
\en{bootstrap} & إعادة سحب الرميات مع الإرجاع لمعرفة كم تتغير النتيجة لو كانت الرميات غيرها\\
\en{covariance} $P$ & جدول عدم يقين المرشح (التباينات) والروابط بينها\\
\en{EKF} & مرشح كالمان الموسَّع: مرشح كالمان لنموذج غير خطي باستخدام جاكوبي\\
\en{family} & مجموعة مقدِّرات تشترك في إعداد ضوضاء مضبوط واحد\\
\en{innovation} & الابتكار: الموقع المقاس ناقص الموقع المتوقع\\
\en{Jacobian} $F$ & مصفوفة المشتقات: كم يتغير كل مُخرَج لخطوة النموذج حين يتغير مُدخَل قليلاً\\
\en{matched / mismatch} & محاكٍ يطابق نموذج المقدِّر / محاكٍ بظواهر واقعية لا يعرفها المقدِّر\\
\en{inverse crime} & الجريمة العكسية: اختبار طريقة على بيانات وُلِّدت بنموذجها نفسه\\
\en{oracle} & العرّاف: مرشح يُعطى الثوابت الحقيقية، حدّ أعلى للمرشحات ذات الثوابت الثابتة\\
\en{paired} & مقترن: مقارنة على الرميات نفسها وبالضوضاء نفسها\\
\en{p-value} & احتمال فرق بهذا الحجم على الأقل لو كانت الطريقتان متساويتين فعلاً\\
$q$، $\sigma$ & ضوضاء العملية (الثقة بالنموذج) وضوضاء القياس (الثقة بالكاميرا)\\
\en{RMS} & الجذر التربيعي لمتوسط مربعات الأخطاء، يعاقب الأخطاء الكبيرة أكثر\\
\en{StaleGuard} & يعتبر القياس المكرَّر بتطابق تام مفقوداً\\
\en{test (dev)} & الاختبار التطويري: مجموعة الاختبار القديمة التي نُظر إليها، للتطوير فقط\\
\en{grid edge} & حافة الشبكة: أفضل قيمة هي الأصغر أو الأكبر المجرَّبة، فقد يكون الأفضل خارج الشبكة\\\bottomrule
\end{tabularx}\end{center}

\end{document}
